\clearpage
\section{Apendice: mapa pagina por pagina de los slides dinamicos}
Este apendice lista cada pagina/slide de la carpeta fuente y su contenido principal. Algunas paginas son portadas, separadores o cierres; se conservan para trazabilidad.
\subsection*{CS3014\_\_\_Heavy\_Light\_Decomposition.pdf (6 paginas)}
\begin{description}[leftmargin=!,labelwidth=4.2em,style=nextline]
\item[Pag. 1] CS3014 - Estructuras de Datos Avanzadas Notas de clase 01 Dynamic Graphs - Heavy-light Decomposition 8 de junio de 2026 Prof. V' ictor Racs' o Galv' an Oyola 1. Descomposiciones de ' arboles Hasta ahora hemos tenido algunas ideas para poder realizar consultas espec' ificas (Euler Tour para LCA) o consultas/modificaciones en sub' arboles (DFS Timestamps) de manera eficiente. Sin embargo, tambi' en es posible recibir consultas/modificaciones sobre las aristas/nodos del camino entre un par de nodos. Para ello, podemos...
\item[Pag. 2] CS3014 - Estructuras de Datos Avanzadas Notas de clase 01 respectivas. 1.1.1. Mayor que la mitad Arista pesada:Aquella arista (u, v) tal que el tama\textasciitilde{} no del sub' arbol deves mayor que la mitad el tama\textasciitilde{} no del sub' arbol deu. (2subtree[v]> subtree[u]). 1.1.2. Mayor que sus hermanos Arista pesada:Aquella arista (u, v) tal que el tama\textasciitilde{} no del sub' arbol deves mayor que los tama\textasciitilde{} nos de los sub' arboles de sus hermanosw, si hay alg' un empate, se considera solo una como pesada. (subtree[v]>=subtree[w], w in children[u]...
\item[Pag. 3] CS3014 - Estructuras de Datos Avanzadas Notas de clase 01 8\} 9int v = G[u][i]; 10if(v == p) continue; 11DFS(v, u); 12if(subtree[v] > subtree[G[u][0]])\{ 13swap(G[u][i], G[u][0]); 14\} 15subtree[u] += subtree[v]; 16\} 17if(p != -1)\{ 18G[u].pop\_back(); 19\} 20\} 21 22void HLD(int u)\{ 23in[u] = T++; 24for(int v : G[u])\{ 25par[v] = u; 26nxt[v] = (v == G[u][0] ? nxt[u] : v); 27h[v] = h[u] + 1; 28HLD(v); 29\} 30out[u] = T - 1; 31\} Notemos que estaremos guardando algunos arreglos extra: inyoutson los timestamps de entrada y sal...
\item[Pag. 4] CS3014 - Estructuras de Datos Avanzadas Notas de clase 01 1int path(int a, int b)\{ 2int res = -2e9; 3while(nxt[a] != nxt[b])\{ 4if(h[nxt[a]] < h[nxt[b]]) swap(a, b); 5res = max(res, query(in[nxt[a]], in[a])); 6a = par[nxt[a]]; 7\} 8if(h[a] > h[b]) swap(a, b); 9res = max(res, query(in[a], in[b])); 10return res; 11\} Es sencillo notar que el bucle se ejecutar' a hasta que est' en ambos nodos en la misma cadena, para lo cual realizamos una consulta extra para cubrir dicho rango de nodos. La complejidad de este algoritmo ...
\item[Pag. 5] CS3014 - Estructuras de Datos Avanzadas Notas de clase 01 Consideremos ahora el caso en el que se nos da un ' arbolT, sobre el cual nosotros queremos responder a cierta informaci' on de cada uno de sus nodosu, considerando que dicha informaci' on depende del sub' arbol deu. Es importante recordar la definici' on de arista pesada y ligera, pues debido a esta se da que existen una cantidad deO(logn) aristas ligeras en el camino desde la ra' iz del ' arbol hasta cualquier nodo u. Entonces, si pudi' eramos procesar cad...
\item[Pag. 6] CS3014 - Estructuras de Datos Avanzadas Notas de clase 01 7int v = G[u][i]; 8if(v == p) continue; 9DFS(v, u); 10if(subtree[v] > subtree[G[u][0]])\{ 11swap(G[u][i], G[u][0]); 12\} 13subtree[u] += subtree[v]; 14\} 15if(p != -1)\{ 16G[u].pop\_back(); 17\} 18\} 19 20void add(int u, int val, int start)\{ 21if(val == 1) agregarInfo(u); 22else quitarInfo(u); 23for(int i = start; i < G[u].size(); i++)\{ 24add(G[u][i], val, 0); // Solo debemos restringir en el primer nivel 25\} 26\} 27 28void solve(int u, int keep)\{ 29for(int i = 1; i...
\end{description}
\subsection*{eda\_slides\_08\_dynamic\_graph.pdf (19 paginas)}
\begin{description}[leftmargin=!,labelwidth=4.2em,style=nextline]
\item[Pag. 1] Dynamic Graph CS3014 - Estructura de Datos Avanzados Luciano A. Romero Calla lromeroc@utec.edu.pe 2026-1
\item[Pag. 2] Overview 1. Warming Up 2. The Dynamic Tree Problem 3. Structural Tool: Preferred Paths 4. Heavy-Light Decomposition (HLD) 5. Summary 2 / 16
\item[Pag. 3] Warming Up 1. Warming Up 3 / 16
\item[Pag. 4] Warming Up Warming Up Problem Remember the problem of Imperial Roads. Now the king may want to keep a chosen road. How does this change your previous solution? A B C D E F 50 70 80 95 20 30 40 50 60 4 / 16
\item[Pag. 5] The Dynamic T ree Problem 2. The Dynamic T ree Problem 5 / 16
\item[Pag. 6] The Dynamic T ree Problem The Dynamic T ree Problem The Challenge:Maintain a forest of rooted trees under: - Structural updates (Link/Cut). - Connectivity queries (FindRoot). - Path aggregates (Min/Max/Sum). Efficiency Goal O(logn)amortized time per operation. r v w link(v, w)? 6 / 16
\item[Pag. 7] Structural T ool: Preferred Paths 3. Structural T ool: Preferred Paths 3.1 Internal Representation: Auxiliary Trees 3.2 The Foundation:access(v) 7 / 16
\item[Pag. 8] Structural T ool: Preferred Paths Structural T ool: Preferred Paths We decompose the tree into vertex-disjoint paths based on access patterns. - Preferred Child:The childcofvsuch that the last access inv's subtree was inc's subtree. - Preferred Edge:Connection to preferred child. - Preferred Path:Maximal chain of preferred edges. 1 2 3 4 5 6 Preferred Path Normal Edge 8 / 16
\item[Pag. 9] Structural T ool: Preferred Paths Internal Representation: Auxiliary T rees Internal Representation: Auxiliary T rees Each preferred path is stored in an Auxiliary Tree (Splay Tree). - Order:Keyed bydepthin the original tree. - Left Subtree:Ancestors (smaller depth). - Right Subtree:Descendants (larger depth). 2 1 4 Splay for Path\{1,2,4\} 3 6 path-parent 9 / 16
\item[Pag. 10] Structural T ool: Preferred Paths The Foundation:access(v) The Foundation:access(v) This operation forces the root-to- v path to become a single preferred path. 1. Splayvto the root of its Aux-Tree. 2. Discardv's right child (it becomes a separate path). 3. Climbvia path-parent tow=pp(v). 4. Splaywand makevits newright child. 5. Repeatuntil the represented root is reached. Effect Afteraccess(v), nodevis at the root of a Splay tree containing exactly the path fromvto the tree root. 10 / 16
\item[Pag. 11] Heavy-Light Decomposition (HLD) 4. Heavy-Light Decomposition (HLD) 4.1 HLD: Why the Logarithmic Bound? 4.2 Complexity Proof Sketch 11 / 16
\item[Pag. 12] Heavy-Light Decomposition (HLD) Heavy-Light Decomposition (HLD) To proveO(logn), we use HLD as an analysis tool (not part of the implementation). Definition Let size(v)be the number of nodes inv's subtree. - Edge(u,v)isHeavyif size(v)> 1 2 size(u). - Edge(u,v)isLightotherwise. Crucial Lemma:Any path from root to node contains at mostlog 2 nlight edges. 12 / 16
\item[Pag. 13] Heavy-Light Decomposition (HLD) HLD: Why the Logarithmic Bound? HLD: Why the Logarithmic Bound? Moving down aLight Edgereduces the subtree size by at least half. 16 9 6 4 4 5 H L (6<=16/2) Max Light Edges\textasciitilde{}log 2(16) =4 13 / 16
\item[Pag. 14] Heavy-Light Decomposition (HLD) Complexity Proof Sketch Complexity Proof Sketch The cost of access(v) is the number of preferred- edge changes. Preferred-Child Changes - Light Changes:Accessing a light child. Onlylognexists on any path. - Heavy Changes:Accessing a heavy child. Can be many, but theycreatelight edges for future operations. Total Cost By potential function analysis (based on HLD), the number of heavy-to-light transitions is bounded. Total:O(logn)amortized. 14 / 16
\item[Pag. 15] Summary 5. Summary 15 / 16
\item[Pag. 16] Summary Summary - Link-Cut Treessolve dynamic forest problems inO(logn). - Key Subroutine:access(v)handles the path logic. - Analysis:Heavy-Light Decomposition provides the amortized guarantee. Operation Mechanism Complexity Access(v) Climbing Aux TreesO(logn) Link(v, w) Access + pp linkO(logn) Cut(v) Access + remove leftO(logn) PathSum(v) Access + Splay AggregateO(logn) 16 / 16
\item[Pag. 17] Thanks
\item[Pag. 18] References References 18 / 16
\item[Pag. 19] Acknowledgements Acknowledgements The course slides are based on the lectures from previous editions and on similar lectures elsewhere. List of credits: Erick Demaine, Keith Schwarz 19 / 16
\end{description}
\subsection*{eda\_slides\_09\_dynamic\_graph.pdf (14 paginas)}
\begin{description}[leftmargin=!,labelwidth=4.2em,style=nextline]
\item[Pag. 1] Dynamic Graph CS3014 - Estructura de Datos Avanzados Luciano A. Romero Calla lromeroc@utec.edu.pe 2026-1
\item[Pag. 2] Overview 1. Dynamic Connectivity 2. Euler-Tour Trees 2 / 11
\item[Pag. 3] Dynamic Connectivity 1. Dynamic Connectivity 1.1 Problem Categorization 1.2 Tree-Structured Solutions: A Contrast 3 / 11
\item[Pag. 4] Dynamic Connectivity Dynamic Connectivity Goal Maintain an undirected graph subject to updates and queries. Operations - insert(u, v): Add edge(u,v). - delete(u, v): Remove edge(u,v). - connected(u, v): Areuandvin the same connected component? u w v insert(u,v) 4 / 11
\item[Pag. 5] Dynamic Connectivity Problem Categorization Problem Categorization Fully Dynamic Admits arbitrary interleavings of edge insertions and deletions. Incremental / Decremental Partially dynamic variants restriction toonlyinsertions oronlydeletions. 5 / 11
\item[Pag. 6] Dynamic Connectivity T ree-Structured Solutions: A Contrast T ree-Structured Solutions: A Contrast Before considering general graphs, we resolve dynamic connectivity withinforests. Metric / Feature Link-Cut Trees (Sleator-Tarjan) Euler-Tour Trees (Henzinger-King) Primary LayoutHeavy-Light Decomposition Flattened Traversal (Euler Tour) Aggregate BoundsPath Aggregates (O(logn)) Subtree Aggregates (O(logn)) ImplementationComplex splay-tree interlinking Single balanced BST representation Edge QueriesDifficult to locate...
\item[Pag. 7] Euler-T our T rees 2. Euler-T our T rees 7 / 11
\item[Pag. 8] Euler-T our T rees Euler-T our T ree Mechanics Concept:Convert a treeTinto a sequential format by tracking its traversal boundary. Every undirected edge is spanned exactly twice (downward and upward). - Each node visit maps to an index in a sequence. - Sequence backed by a balanced BST (e.g., Treap, AVL). - Keys are implicitly defined by sequential array order. 1 2 3 Sequence:1->2->2->1->3->3->1 8 / 11
\item[Pag. 9] Euler-T our T rees Structural Mutations: Link \& Cut By tracking pointers to the first and last occurrence of nodev, we manipulate subtrees inO(logn) time via basic BST splits and concatenations. Cut(v) Removes the subtree rooted atv. 1.Identifyv's first and last visit in the BST. 2.Split sequence into three:Before-v, Subtree-v,After-v. 3.ConcatenateBefore-vandAfter-v. Link(u, v) Attachesu's tree as a child ofv. 1.Splitv's sequence right before its final occurrence. 2.Splice the entire sequence ofuinto the gap. 3.Me...
\item[Pag. 10] Euler-T our T rees Fully Dynamic Connectivity in General Graphs How do we handle general graphs rather than isolated trees? - Maintain spanning forests of the graph. - Partition edges intolognlevels based on assigned "weights" or frequencies. - Use an Euler-Tour tree to maintain the spanning forest of each level. - Amortized Cost:O(log 2 n)for updates,O(logn)for queries. LeveliForest (BST) Leveli+1 Forest edge push 10 / 11
\item[Pag. 11] Euler-T our T rees Summary \& Lower Bounds Key Takeaways: - Euler-tour trees are simpler to implement and analyze than Link-Cut trees for aggregate subtree information. - General graph dynamic connectivity can be achieved inO(log2 n)time using level-partitioned Euler-Tour trees. Lower Bounds (P atrascu \& Demaine): - O(xlogn)update requires (logn/logx)query time. - O(xlogn)query requires (lognlogx)update time. - Conclusion:There is an inherent trade-off in dynamic graph structures! 11 / 11
\item[Pag. 12] Thanks
\item[Pag. 13] References References 13 / 11
\item[Pag. 14] Acknowledgements Acknowledgements The course slides are based on the lectures from previous editions and on similar lectures elsewhere. List of credits: Erick Demaine, Keith Schwarz 14 / 11
\end{description}
